\subsection{A graph-universal scalar homotopy and bounded cyclic blocks}
\label{subsec:graph-homotopy}

Let $\bar o:=V\setminus\{o\}$,
$H:=Q_{\bar o\bar o}$, and
$a:=-Q_{\bar o,o}\geq0$.  Fixing the root coordinate at $t\geq0$ gives the
conditional obstacle problem
\begin{equation}
 z(t)
 :=\argmin_{z\in\mathbb R_+^{\bar o}}
 \left\{
  \frac12z^\top Hz-(c_{\rho,\bar o}+at)^\top z
 \right\},
 \qquad
 \mu(t):=a^\top z(t).
 \label{eq:graph-conditional-response}
\end{equation}

\begin{theorem}[Graph-universal scalar active-set homotopy]
\label{thm:graph-scalar-homotopy}
For every finite graph with one seed:
\begin{enumerate}[label=(\roman*),leftmargin=*]
 \item $z(t)$ is continuous and coordinatewise nondecreasing; its supports
 are nested and change at most $n-1$ times.
 \item On an interval with active set $A\subseteq\bar o$,
 \begin{align}
  z_A(t)&=H_{AA}^{-1}(c_{\rho,A}+a_At),
  \label{eq:graph-homotopy-affine}\\
  \mu'(t)&=a_A^\top H_{AA}^{-1}a_A,
  \label{eq:graph-homotopy-response-slope}\\
  G'(t)&=Q_{oo}-a_A^\top H_{AA}^{-1}a_A\geq\alpha,
  \qquad
  G(t):=Q_{oo}t-c_{\rho,o}-\mu(t).
  \label{eq:graph-root-schur-slope}
 \end{align}
 \item If $c_{\rho,o}>0$, the root value of $x^\star(\rho)$ is the unique
 zero of $G$.  Walking the affine pieces of $G$ discovers the exact support
 without a support oracle.  If $c_{\rho,o}\leq0$, the optimum is zero.
\end{enumerate}
\end{theorem}

\begin{proof}
Strong convexity gives uniqueness and continuity.  Increasing $t$ increases
the demand by the nonnegative vector $a(t_2-t_1)$, so Stieltjes obstacle
comparison gives $z(t_2)\geq z(t_1)$.  Thus supports are nested and each
change activates at least one new vertex.  On a fixed support, the active KKT
equalities give \cref{eq:graph-homotopy-affine} and differentiation gives
\cref{eq:graph-homotopy-response-slope}.  The final expression in
\cref{eq:graph-root-schur-slope} is the Schur complement of $H_{AA}$ in the
principal matrix on $A\cup\{o\}$.  Since that matrix is bounded below by
$\alpha I$, its scalar Schur complement is at least $\alpha$.

For completeness, the next event is locally explicit.  At an inactive
coordinate $v$, define its obstacle slack
\begin{equation}
 g_v(t)
 :=H_{vA}H_{AA}^{-1}(c_{\rho,A}+a_At)
   -c_{\rho,v}-a_vt.
 \label{eq:graph-homotopy-slack}
\end{equation}
Its slope is nonpositive because $H_{vA}\leq0$,
$H_{AA}^{-1}a_A\geq0$, and $a_v\geq0$.  Only a graph-boundary vertex can
have negative slope.  Starting from $A=\emptyset$ at $t=0$, take the smallest
future zero of the affine slacks and add every tied coordinate.  Between two
such events, $G$ is affine.  If its affine zero occurs no later than the next
event, stop; otherwise consume the event and continue.  Strictly negative
nonseed demands prevent a nonboundary coordinate from appearing in the same
event without first becoming exposed.  Hence the walk is adjacency-local and
consumes precisely the labels positive at the final root value.  The active
root equation is $G(t)=0$, and \cref{eq:graph-root-schur-slope} makes its zero
unique.  When $c_{\rho,o}\leq0$, zero satisfies all shifted KKT inequalities.
\end{proof}

\begin{lemma}[One-event Schur update]
\label{lem:one-event-schur-update}
Suppose the homotopy has active set $A$ immediately before an untied event
$k\notin A$ at $t=\tau$.  Put
\begin{equation}
 M:=H_{AA},
 \qquad h:=H_{Ak},
 \qquad r:=M^{-1}h\leq0,
 \qquad \sigma:=H_{kk}-h^\top M^{-1}h>0.
 \label{eq:event-schur-data}
\end{equation}
On the next affine piece,
\begin{align}
 z_k^+(t)&=-\frac{g_k(t)}{\sigma},
 &z_A^+(t)&=z_A(t)-r z_k^+(t),
 \label{eq:event-primal-update}\\
 g_v^+(t)&=g_v(t)+s_{vk}z_k^+(t),
 &s_{vk}&:=H_{vk}-H_{vA}M^{-1}H_{Ak}\leq0
 \qquad(v\notin A\cup\{k\}).
 \label{eq:event-slack-update}
\end{align}
In particular, since $g_k(\tau)=0$ and $g_k'(t)<0$ at a genuine untied
crossing, every remaining slack slope can only become more negative.  Thus
one activation changes all future event keys through one column of a Schur
complement.  A tied event has the identical block-Schur formula.
\end{lemma}

\begin{proof}
The old active equations are $Mz_A=d_A(t)$.  After adjoining $k$, eliminate
the old coordinates from
\begin{equation*}
 \begin{bmatrix}M&h\\h^\top&H_{kk}\end{bmatrix}
 \begin{bmatrix}z_A^+\\z_k^+\end{bmatrix}
 =\begin{bmatrix}d_A(t)\\d_k(t)\end{bmatrix}.
\end{equation*}
The first block row gives $z_A^+=z_A-rz_k^+$; substitution into the second
gives $\sigma z_k^+=-g_k(t)$.  Insert this expression in the definition of
an inactive slack to obtain \cref{eq:event-slack-update}.  Principal Schur
complements of a Stieltjes matrix are again Stieltjes, so
$\sigma>0$, $r\leq0$, and $s_{vk}\leq0$ for $v\ne k$.  The slope claim
follows from $z_k^+(t)=-g_k'(t)(t-\tau)/\sigma$.
\end{proof}

\begin{lemma}[Exact tied-batch Schur update]
\label{lem:block-event-schur-update}
Let $A$ be the current active set and let a nonempty block
$B\subseteq\bar o\setminus A$ join on the next homotopy piece.  Write
\begin{equation}
 M:=H_{AA},
 \qquad E:=H_{AB},
 \qquad
 \mathcal S_B:=H_{BB}-E^\top M^{-1}E,
 \label{eq:block-event-schur-data}
\end{equation}
and let
$g_B(t):=H_{BA}z_A(t)-d_B(t)$, where
$d(t):=c_{\rho,\bar o}+at$.  Then the new affine solution satisfies
\begin{align}
 z_B^+(t)&=-\mathcal S_B^{-1}g_B(t),
 &z_A^+(t)&=z_A(t)-M^{-1}Ez_B^+(t),
 \label{eq:block-event-primal-update}\\
 g_C^+(t)&=g_C(t)+\mathcal S_{CB}z_B^+(t),
 &\mathcal S_{CB}&:=H_{CB}-H_{CA}M^{-1}E\leq0
 \label{eq:block-event-slack-update}
\end{align}
for every remaining inactive block
$C\subseteq\bar o\setminus(A\cup B)$.  Moreover,
$\mathcal S_B$ is Stieltjes, $\mathcal S_B^{-1}\geq0$, and
$-M^{-1}E\geq0$.  The same identities hold for an exact boundary-gate batch
after replacing $H$ by $Q$ and $g_B$ by the current shifted KKT slack.
\end{lemma}

\begin{proof}
The enlarged active equations are
\begin{equation*}
 \begin{bmatrix}M&E\\E^\top&H_{BB}\end{bmatrix}
 \begin{bmatrix}z_A^+\\z_B^+\end{bmatrix}
 =\begin{bmatrix}d_A(t)\\d_B(t)\end{bmatrix}.
\end{equation*}
Eliminating the old coordinates gives
$\mathcal S_Bz_B^+=d_B-H_{BA}z_A=-g_B$ and then the second formula in
\cref{eq:block-event-primal-update}.  Substitution in the definition of the
remaining inactive slacks gives
\cref{eq:block-event-slack-update}.  Principal Schur complements of a
Stieltjes matrix are Stieltjes, and inverse positivity gives all displayed
signs.  The gate statement is the identical block elimination applied to
$Q_{A\cup B,A\cup B}$.
\end{proof}

The theorem removes cyclic ambiguity at the level of correctness: there is
still one ordered stream of activation events.  A literal implementation,
however, refactors the current principal system and refreshes every boundary
slack after each event, which can cost quadratic output work.
\Cref{lem:one-event-schur-update} shows exactly what must be compressed: a
kinetic minimum under successive signed Schur-column or block-column updates;
\cref{lem:block-event-schur-update} shows that simultaneous crossings require
no additional algebra.  Trees avoid
that refresh through scalar child independence.  The same mechanism extends
to cyclic graphs whose articulation blocks have bounded size.

Count every bridge as a two-vertex block.  Suppose the block-cut
decomposition is exposed with the adjacency oracle, and every biconnected
block contains at most $q$ vertices.  This representation identifies blocks
but supplies neither the RPPR support nor its radius.

\begin{theorem}[Condition-free product bound for bounded cyclic blocks]
\label{thm:bounded-block-response}
Under the preceding block-incidence representation, exact single-seed RPPR
can be solved using
\begin{align}
 \Work_{\rm block}
 &=O\!\left(
  q^3\log(1+\Delta)
  \left[
   \vol(S^\star)
   +\sum_{v\in S^\star}\operatorname{dist}(o,v)
  \right]
 \right)\notag\\
 &=\widetilde O\!\left(
  \frac{q^3}{\rho\sqrt\alpha}
 \right)
 \label{eq:bounded-block-work}
\end{align}
degree work and $O(q\vol(S^\star))$ storage, without a support, radius, or
confinement oracle.  Thus every fixed-$q$ block graph, including graphs formed
from trees by gluing bounded-size cyclic gadgets at articulation vertices,
meets the intended product scale.
\end{theorem}

\begin{proof}
Root the block-cut tree at $o$.  A descendant block $K$ meets its parent side
in one articulation $p$, so after fixing $x_p=t$ the entire block-descendant
region sends a scalar response to $p$.  Its first event is obtained by scanning
the at most $q$ vertices of $K$ adjacent to $p$; no descendant block of an
inactive articulation can activate first.

Maintain one lazy iterator per exposed block.  On a fixed collection of child
response pieces and a fixed active subset of $K$, the at most $q-1$ internal
coordinates solve one affine linear system in $t$.  The next event is either
an inactive internal KKT slack reaching zero or an active articulation value
reaching the next breakpoint of one of its child-block responses.  Both are
zeros of known affine functions.  Refactor the local system and refresh its
event heap in $O(q^3\log(1+\Delta))$ work.  Stieltjes comparison and
\cref{thm:graph-scalar-homotopy} guarantee that coordinates never deactivate
and that these are all possible events.

When an event labeled by $v$ is consumed, advance it through each ancestor
block on the block-cut path to the root.  Distinct consecutive blocks contain
distinct articulation vertices joined by a nonempty graph path, so the number
of traversed blocks is at most $1+\operatorname{dist}(o,v)$.  Exposing a block
scans only a bounded-size incident gadget; summing these scans over active
articulations and their inactive boundary blocks is
$O(q\vol(S^\star))$.  The root iterator stops at the unique zero in
\cref{thm:graph-scalar-homotopy}; its consumed labels are exactly
$S^\star$.  One final width-$(q-1)$ block solve and boundary KKT scan recovers
and certifies $x^\star$ within the displayed order.  Apply
\cref{lem:graph-support-radius} and
$\vol(S^\star)\leq1/\rho$ to obtain the second line of
\cref{eq:bounded-block-work}.
\end{proof}

The preceding representation is convenient but not necessary.  The missing
fact is again causal.  An inactive vertex contributes no off-diagonal term to
an active equation.  Consequently a cycle can become relevant only at the
event at which one of its previously inactive vertices has value zero.

\begin{lemma}[Zero-prefix invariance under causal block merges]
\label{lem:causal-block-zero-prefix}
Run the root-conditioned homotopy of
\cref{thm:graph-scalar-homotopy} using ordinary degree queries and adjacency
scans, and scan a vertex only when it activates.  Let $A$ be the active set
immediately before a homotopy event at parameter $t_0$, and let $B$ be its
tied activation batch.  The exposed graph contains every cached edge with at
least one scanned endpoint, including its zero-valued inactive boundary
endpoint.  After scanning $B$, insert every newly revealed edge into this
exposed graph.  Then:
\begin{enumerate}[label=(\roman*),leftmargin=*]
 \item every edge newly inserted into the exposed graph has at least one
 endpoint in $B$;
 \item the old solution $z_A(t_0)$, extended by $z_B(t_0)=0$, is the unique
 solution of the enlarged active equations at $t_0$, while every other newly
 exposed endpoint remains a valid zero boundary coordinate;
 \item every response piece already consumed at a parameter below $t_0$
 remains exact.  Only unconsumed lookaheads and future affine pieces belonging
 to blocks changed by the insertion can become invalid.
\end{enumerate}
The same statements hold when several edges inside $B$ merge more than two
current blocks at once.
\end{lemma}

\begin{proof}
Every newly revealed edge is read from an adjacency list in $B$, proving
(i); an edge with two previously scanned endpoints was already cached.  Before
$t_0$ every coordinate in $B$ is zero.  Hence an incidence from $B$ to $A$,
an incidence internal to $B$, and an incidence from $B$ to a future boundary
vertex contribute zero to the old active equations.  At $t_0$ the
KKT slack of every label in the tied batch is zero by definition.  Therefore
the old equations on $A$ and the new equations on $B$ are all satisfied by
$(z_A(t_0),0_B)$.  The enlarged principal matrix is positive definite, so
this solution is unique, proving (ii).

For every earlier parameter, the same zero-coordinate argument shows that
the response seen by the old active variables is unchanged.  Continuity at
$t_0$ gives a common checkpoint for the old and enlarged systems.  Refactoring
from that checkpoint changes their derivatives and future breakpoints but
not their common prefix, proving (iii).  Internal batch edges have two zero
endpoints, so the proof is unchanged for a multiway merge.
\end{proof}

\begin{theorem}[Representation-free product bound for bounded cyclic blocks]
\label{thm:causal-bounded-block-response}
Suppose every biconnected block in the seed component has at most $q$
vertices, with each bridge counted as a two-vertex block.  The block
decomposition, block identifiers, root orientation, and $q$ itself need not
be supplied.  With only ordinary degree queries and adjacency-list scans,
exact single-seed RPPR can be solved in
\begin{align}
 \Work_{\rm causal\text{-}block}
 &=O\!\left(
   q^3\log\!\bigl(2+\Delta+\vol(S^\star)\bigr)
   \sum_{v\in S^\star}
    (1+d_v)\bigl(1+\operatorname{dist}(o,v)\bigr)
  \right)\notag\\
 &=\widetilde O\!\left(
   \frac{q^3}{\rho\sqrt\alpha}
  \right)
 \label{eq:causal-bounded-block-work}
\end{align}
degree work and $O(q\vol(S^\star))$ storage, without a support, radius,
confinement, or structural oracle.  In particular, the presentation
condition in \cref{thm:bounded-block-response} can be removed without losing
the intended product scale.
\end{theorem}

\begin{proof}
Maintain the block--cut forest of the exposed graph: every scanned vertex,
its inactive boundary labels, and every cached incidence are present, while
unscanned boundary coordinates are fixed at zero.  This convention is
important.  If one boundary label touches several active branches, its
combined slack already belongs to one exposed block response rather than to a
separate cross-branch demand array.  Adding a newly scanned adjacency list
consists of leaf insertions and contractions of block--cut paths: inserting a
cycle-closing edge merges precisely the blocks and articulations on the path
between its endpoints.  For a tied batch, inserting all of its edges
contracts the corresponding union of such paths.  Every current biconnected
block is contained in a biconnected block of the full graph, so every
contracted path or subtree that becomes one block contains at most $q$
original vertices.  A rooted dynamic-forest path expose, followed by
union--find contraction, therefore maintains the combinatorial block forest
in
\begin{equation}
 O\!\left(q\log\bigl(2+\vol(S^\star)\bigr)\right)
 \label{eq:causal-block-forest-update}
\end{equation}
work per inserted exposed incidence.  Processing all incidences of a tied
batch before publishing its new block records avoids any dependence on their
arbitrary scan order.

Keep the lazy scalar block iterators of
\cref{thm:bounded-block-response}, but make every response record checkpointed
and versioned, with one current piece and one live lookahead.  At an
articulation, its off-path child iterators remain
behind one meldable-heap handle owned by the articulation rather than by its
current parent block.  Thus a contraction does not enumerate or copy those
child subtrees.  Blocks disjoint from the contracted subtree retain their
iterators.  For every changed block, discard stale unconsumed heap records,
checkpoint all incident child iterators at their current articulation values,
and rebuild the at-most-$q$-variable affine system and its next-event heap.
Dense refactorization and rekeying cost
\begin{equation}
 O\!\left(q^3\log\bigl(2+\Delta+\vol(S^\star)\bigr)\right)
 \label{eq:causal-block-rebuild}
\end{equation}
per conservatively charged newly exposed incidence.  No past breakpoint list is
replayed.  By \cref{lem:causal-block-zero-prefix}, the checkpoint and every
consumed prefix are exact; the rebuilt iterator is the unique exact future
continuation.  Induction over activation batches therefore proves exactness.

Every admitted adjacency list is scanned once.  Charge its exposed-incidence
insertions, stale versions, local rebuilds, and any conservative propagated
rekey to the newly activated endpoint $v$.  A propagated version crosses at
most $1+\operatorname{dist}(o,v)$ block nodes, just as in
\cref{thm:bounded-block-response}.  Monotonicity labels every genuinely
consumed response event by one support vertex, and equality remains inactive,
so there are no additional unlabeled event chains.  Summing the per-incidence
charge gives the first line of \cref{eq:causal-bounded-block-work}; final
block-tree recovery and the cached boundary KKT scan fit in the same bound.
Addressable handles replace invalidated ancestor keys in place, while a
consumed prefix retains only its emitted labels and current checkpoints.
Consequently current dense factors, articulation heaps, cached incidences,
and version records use $O(q\vol(S^\star))$ storage.

If the seed component is nontrivial, $1+d_v\leq2d_v$ for every support
vertex.  Writing $R_\star:=\max_{v\in S^\star}\operatorname{dist}(o,v)$,
\begin{equation*}
 \sum_{v\in S^\star}(1+d_v)
       \bigl(1+\operatorname{dist}(o,v)\bigr)
 \leq2(1+R_\star)\vol(S^\star).
\end{equation*}
The isolated-seed case is constant work.  Now apply
\cref{lem:graph-support-radius} and
$\vol(S^\star)\leq1/\rho$ to obtain the second line of
\cref{eq:causal-bounded-block-work}.  The algorithm uses realized block sizes
only; $q$ is an analysis parameter, not advice.
\end{proof}

Large biconnected cores remain different, but block size alone is not the
right obstruction.  \Cref{thm:cycle-one-pass} closes an arbitrarily long
cycle in output-linear work, \cref{thm:equitable-shell-one-pass} closes
arbitrarily thick root-distance-equitable shells in the same work,
\cref{thm:hidden-equitable-tree-quotient} closes multiple thick cell types
whenever an equitable tree quotient exists, without supplied cell labels, and
\cref{cor:radial-width-gate} below closes every bounded-width core whose true
support is thin across distance shells.  The open data-structure lemma is now
more precise: in a core with neither a bounded-block decomposition, an
equitable tree quotient, nor thin radial support, compress the scalar homotopy
trace without refreshing all boundary slacks after every activation, or prove
that the exact boundary gate realizes only
$\widetilde O(1/\sqrt\alpha)$ batches.  The block update in
\cref{lem:block-event-schur-update} is the exact algebraic interface for the
first route; \cref{prop:three-path-gap-obstruction} below shows that the second
route cannot follow from a uniform accelerated objective contraction.
